\subsection{$\epsilon-\delta$ Limits, Uniform Continuity, and the Bolzano-Weierstrass Theorem}

Building upon the topological framework of $\mathbb{R}$, we rigorously define the limiting behavior of functions and sequences, transitioning from local properties to global guarantees on compact domains.

\begin{definition}
Let $E \subset \mathbb{R}$, and let $c$ be a limit point of $E$. A function $f: E \to \mathbb{R}$ has a \textbf{limit} $L \in \mathbb{R}$ as $x$ approaches $c$, denoted $\lim_{x \to c} f(x) = L$, if for every $\epsilon > 0$, there exists a $\delta > 0$ such that for all $x \in E$,
\[
0 < |x - c| < \delta \implies |f(x) - L| < \epsilon
\]
\end{definition}

From this local definition, we define continuity: a function $f$ is continuous at $c \in E$ if $\lim_{x \to c} f(x) = f(c)$. If this holds for all points in $E$, $f$ is continuous on $E$. While continuity is a local property (the choice of $\delta$ may depend on both $\epsilon$ and $c$), uniform continuity establishes a global bound independent of $c$.

\begin{definition}
A function $f: E \to \mathbb{R}$ is \textbf{uniformly continuous} on $E$ if for every $\epsilon > 0$, there exists a $\delta > 0$ such that for all $x, y \in E$,
\[
|x - y| < \delta \implies |f(x) - f(y)| < \epsilon
\]
\end{definition}

\begin{theorem}[Heine-Cantor Theorem]
If $K \subset \mathbb{R}$ is compact and $f: K \to \mathbb{R}$ is continuous, then $f$ is uniformly continuous on $K$.
\end{theorem}

\begin{proof}
Let $\epsilon > 0$ be given. Because $f$ is continuous on $K$, for each $p \in K$, there exists a $\delta_p > 0$ such that for all $x \in K$,
\[
|x - p| < \delta_p \implies |f(x) - f(p)| < \frac{\epsilon}{2}
\]
Consider the collection of open neighborhoods $V_p = (p - \frac{\delta_p}{2}, p + \frac{\delta_p}{2})$ for all $p \in K$. The collection $\{V_p\}_{p \in K}$ forms an open cover of $K$. Because $K$ is compact, there exists a finite subcover corresponding to points $p_1, p_2, \dots, p_n \in K$. 

Define $\delta = \frac{1}{2} \min\{\delta_{p_1}, \delta_{p_2}, \dots, \delta_{p_n}\}$. Since the set of radii is finite and each $\delta_{p_i} > 0$, we have $\delta > 0$. Let $x, y \in K$ such that $|x - y| < \delta$. Because the finite collection $\{V_{p_i}\}_{i=1}^n$ covers $K$, there is some integer $k \in \{1, \dots, n\}$ such that $x \in V_{p_k}$. This implies:
\[
|x - p_k| < \frac{\delta_{p_k}}{2}
\]
By the triangle inequality, we bound the distance between $y$ and $p_k$:
\[
|y - p_k| \le |y - x| + |x - p_k| < \delta + \frac{\delta_{p_k}}{2}
\]
Since $\delta \le \frac{\delta_{p_k}}{2}$, we deduce $|y - p_k| < \delta_{p_k}$. We evaluate the difference in functional values using the continuity of $f$ at $p_k$. Both $x$ and $y$ are strictly within $\delta_{p_k}$ of $p_k$, yielding:
\[
|f(x) - f(y)| \le |f(x) - f(p_k)| + |f(p_k) - f(y)| < \frac{\epsilon}{2} + \frac{\epsilon}{2} = \epsilon
\]
Thus, $f$ is uniformly continuous on $K$.
\end{proof}

We conclude this chapter with a fundamental property of sequences in bounded domains, relying heavily on the completeness axiom.

\begin{theorem}[Bolzano-Weierstrass Theorem]
Every bounded sequence in $\mathbb{R}$ possesses a convergent subsequence.
\end{theorem}

\begin{proof}
Let $(x_n)_{n=1}^\infty$ be a bounded sequence in $\mathbb{R}$. There exists a closed interval $I_0 = [a_0, b_0]$ such that $x_n \in I_0$ for all $n \in \mathbb{N}$. We proceed by constructing a nested sequence of closed intervals. 

Bisect $I_0$ into two closed subintervals: $[a_0, \frac{a_0+b_0}{2}]$ and $[\frac{a_0+b_0}{2}, b_0]$. Since $I_0$ contains infinitely many terms of the sequence $(x_n)$, at least one of these two subintervals must contain infinitely many terms. Let $I_1 = [a_1, b_1]$ denote this subinterval. We select an index $n_1$ such that $x_{n_1} \in I_1$.

Assume inductively that we have constructed a closed interval $I_k = [a_k, b_k]$ containing infinitely many terms of $(x_n)$, with length $\frac{b_0-a_0}{2^k}$, and we have chosen an index $n_k > n_{k-1}$ such that $x_{n_k} \in I_k$. Bisect $I_k$ into two equal halves. At least one half contains infinitely many terms of $(x_n)$; call it $I_{k+1} = [a_{k+1}, b_{k+1}]$. Choose an index $n_{k+1} > n_k$ such that $x_{n_{k+1}} \in I_{k+1}$.

This process recursively generates a sequence of nested closed intervals $I_0 \supset I_1 \supset I_2 \supset \dots \supset I_k \supset \dots$ and a subsequence $(x_{n_k})_{k=1}^\infty$ where $x_{n_k} \in I_k$ for all $k$. The sequences of endpoints $(a_k)$ and $(b_k)$ are monotonic: $(a_k)$ is non-decreasing and bounded above by $b_0$, while $(b_k)$ is non-increasing and bounded below by $a_0$. By the Monotone Convergence Theorem (derived from the Least Upper Bound property), both converge. Since the interval length is $b_k - a_k = \frac{b_0-a_0}{2^k}$, we have $\lim_{k \to \infty} (b_k - a_k) = 0$, proving they converge to the identical limit $L \in \mathbb{R}$.

Because $a_k \le x_{n_k} \le b_k$ for all $k$, the Squeeze Theorem forces $\lim_{k \to \infty} x_{n_k} = L$. Thus, $(x_{n_k})$ converges.
\end{proof}